%%%%%%%%%%%%%%%%%%%%%%%%%%%%%%%%%%%%%%%%%%%%%%%%%%%%%%%%%%%%%%%%%%%%%%%%
%% Customizações do abnTeX2 (https://www.abntex2.net.br)              %%
%% para a Empresa Brasileira de Pesquisa Agropecuária (Embrapa)       %%
%%                                                                    %%
%% This work may be distributed and/or modified under the             %%
%% conditions of the LaTeX Project Public License, either version 1.3 %%
%% of this license or (at your option) any later version.             %%
%% The latest version of this license is in                           %%
%%   http://www.latex-project.org/lppl.txt                            %%
%% and version 1.3 or later is part of all distributions of LaTeX     %%
%% version 2005/12/01 or later.                                       %%
%%                                                                    %%
%% This work has the LPPL maintenance status `maintained'.            %%
%%                                                                    %%
%% The Current Maintainer of this work is Igor Lopes                  %%
%% Contact: igor.lopes@embrapa.br                                    %%
%%                                                                    %%
%% Further information about abnTeX2                                  %%
%% are available on https://www.abntex2.net.br/                       %%
%%                                                                    %%
%%%%%%%%%%%%%%%%%%%%%%%%%%%%%%%%%%%%%%%%%%%%%%%%%%%%%%%%%%%%%%%%%%%%%%%%

\documentclass[
    a4paper,          % Tamanho da folha A4
    12pt,             % Tamanho da fonte 12pt
    chapter=TITLE,    % Todos os capitulos devem ter caixa alta
    section=TITLE,    % Todas as secoes devem ter caixa alta
    oneside,          % Usada para impressao em apenas uma face do papel
    english,          % Hifenizacoes em ingles
    spanish,          % Hifenizacoes em espanhol
    brazilian            % Ultimo idioma eh o idioma padrao do documento
]{abntex2}

% Importações de pacotes
\usepackage{babel} 
\usepackage[utf8]{inputenc}                         % Acentuação direta
\usepackage[T1]{fontenc}                            % Codificação da fonte em 8 bits
\usepackage{graphicx}                               % Inserir figuras
\usepackage{amsfonts, amssymb, amsmath}             % Fonte e símbolos matemáticos
\usepackage{booktabs}                               % Comandos para tabelas
\usepackage{verbatim}                               % Texto é interpretado como escrito no documento
\usepackage{multirow, array}                        % Múltiplas linhas e colunas em tabelas
\usepackage{longtable}                              % Tabelas que se estendem por mais de uma página (\EMBRAPAtablonga)
\usepackage{indentfirst}                            % Endenta o primeiro parágrafo de cada seção.
\usepackage{listings}                               % Utilizar codigo fonte no documento
\usepackage{xcolor}
\usepackage[portuguese,ruled,lined]{algorithm2e}    % Escrever algoritmos
\usepackage{amsgen}
\usepackage{tocloft}                                % Permite alterar a formatação do Sumário
\usepackage{etoolbox}                               % Usado para alterar a fonte da Section no Sumário
\usepackage[nogroupskip,nonumberlist,acronym]{glossaries}                % Permite fazer o glossario
\usepackage{caption}                                % Altera o comportamento da tag caption
% Citações e referências ABNT (NBR 10520 / NBR 6023) via biblatex-abnt + biber.
% Substitui o antigo abntex2cite (BibTeX). Mapeamento das opções de estilo:
%   style=abnt        sistema alfabético (autor-data); destaque do título em negrito (padrão)
%   giveninits=true   prenomes abreviados — ex.: OLIVEIRA, C. A. (como no abntex2cite)
%   justify           lista de referências justificada (antigo bibjustif)
%   maxcitenames=3 / mincitenames=1   "et al." nas citações com 4+ autores (antigo abnt-etal-cite=3)
%   maxbibnames=99    lista TODOS os autores em cada referência (antigo abnt-etal-list=0)
% (recuo=0cm e o destaque em negrito já são o padrão do biblatex-abnt; "et al." em
%  itálico, se desejado, é um \DefineBibliographyStrings — ver README.)
\usepackage[style=abnt,backend=biber,giveninits=true,justify,maxcitenames=3,mincitenames=1,maxbibnames=99]{biblatex}
\addbibresource{elementos-pos-textuais/referencias.bib}   % base BibTeX das referências
\usepackage{mathptmx}                               % Usa a fonte Times New Roman
\usepackage{microtype}                              % Melhorias de justificação (após a seleção da fonte)
\usepackage{appendix}                               % Gerar o apendice no final do documento
\usepackage{paracol}                                % Criar paragrafos sem identacao
\usepackage{lib/embrapatex}                         % Biblioteca EmbrapaTex para documentos padronizados
\usepackage{pdfpages}                               % Incluir pdf no documento

% --- Pacotes opcionais (desabilitados) ---
% Descomente conforme a necessidade do documento:
%\usepackage{float}            % floats com posicionamento fixo [H]
%\usepackage[bottom]{footmisc} % fixa as notas de rodapé no rodapé da página
%\usepackage{lscape}           % páginas isoladas em modo paisagem
%\usepackage{subcaption}       % subfiguras/subtabelas (substitui o obsoleto subfig)
%\usepackage{upgreek}          % letras gregas não-itálicas (\upmu, \uppi, ...)

% Organiza e gera a lista de abreviaturas, simbolos e glossario
\makeglossaries

% Gera o Indice do documento
\makeindex
